% Part: lambda-calculus
% Chapter: lambda-definability
% Section: pairs

\documentclass[../../../include/open-logic-section]{subfiles}

\begin{document}

\olfileid{lam}{ldf}{pai}
\olsection{જોડીઓ અને પૂર્વગામી}

\begin{defn}
$M$ અને~$N$ની જોડી ($\tuple{M,N}$ તરીકે લખાતી)
આ રીતે વ્યાખ્યાયિત થાય છે:
\[
\tuple{M,N} \ident \lambd[f][fMN].
\]
\end{defn}
  
સહજ રીતે, આ એવું વિધેય છે જે બીજું વિધેય
સ્વીકારે છે અને તેને જોડીના બંને ઘટકો પર લાગુ
કરે છે. આ વિચાર પરથી બે પદો લઈને તેમની
જોડી આપતું રચનાપદ મળે છે:
\[
  \fn{Pair} \ident \lambd[mn][\lambd[f][fmn]]
\]
જોડી આપેલી હોય ત્યારે તેના ઘટકો પાછા પણ
મેળવવા છે. તેના માટે બે પ્રવેશ-વિધેયો જોઈએ;
તે જોડી આર્ગ્યુમેન્ટ તરીકે લે છે અને અનુક્રમે
પહેલો કે બીજો ઘટક આપે છે:
\begin{align*}
  \fn{Fst} & \ident \lambd[p][p(\lambd[mn][m])]\\
  \fn{Snd} & \ident \lambd[p][p(\lambd[mn][n])]
\end{align*}

\begin{prob}
  પ્રવેશ-વિધેયો $\fn{Fst}$ અને $\fn{Snd}$
  કેમ કામ કરે છે તે સમજાવો.
\end{prob}

હવે જોડીઓ વડે પૂર્વગામી વિધેયને !!{lambda define}
કરી શકીએ:
\[
  \fn{Pred} \ident \lambd[n][\fn{Fst}(n (\lambd[p][\tuple{\fn{Snd}\, {p}, \fn{Succ}(\fn{Snd}\, {p})}]) \tuple{\num 0, \num 0})]
\]
યાદ રાખો કે $\num n\, f x$નું લઘુકરણ
$f^{n}(x)$માં થાય છે. અહીં~$f$ જોડી~$p$
સ્વીકારી જોડી~$p$નો બીજો ઘટક તથા એ બીજા ઘટકનો
અનુગામી ધરાવતી નવી જોડી આપતું વિધેય છે;
$x$ એ જોડી $\tuple{\num 0,\num 0}$ છે.
તેથી $n=0$ માટે પરિણામ $\tuple{\num 0,\num 0}$
અને અન્ય કિસ્સામાં $\tuple{\num{n-1}, \num n}$
મળે છે. પછી~$\fn{Pred}$ એ પરિણામનો
પહેલો ઘટક આપે છે.
\sourcecorrection{OLLAM-052}{સ્થિર અંગ્રેજી મૂળ
ગદ્યમાં શરૂઆતની અને શૂન્ય માટેની જોડીને
$\tuple{0,0}$ લખે છે, જ્યારે પદની વ્યાખ્યામાં
ચર્ચ અંકોની જોડી~$\tuple{\num 0,\num 0}$
છે. અહીં પદ-સ્તરે ચર્ચ અંક જ લખ્યા છે.}

બાદબાકી~$a$ પર $\fn{Pred}$ને $b$ વાર
લાગુ કરીને વ્યાખ્યાયિત કરી શકાય:
\[
\fn{Sub} \ident \lambd[ab][b \fn{Pred}\, a].
\]
    
\end{document}
